\begin{proof}[Proof of \cref{obj:prop:dual-representation}]
Write
\[
p_j=P(X=j),\qquad q_{(a,y)j}=P(X=j,A=a,Y=y),\qquad
\mu_{aj}=E_P[Y(a)\mid X=j],\qquad
\tau_j=\mu_{1j}-\mu_{0j},\qquad
B(P)=\sum_{j=1}^d p_j\mu_{0j}.
\]
For a cell with \(p_j=0\), set \(\mu_{aj}=0\); its contributions below vanish. The cell masses are nonnegative, and the binary outcomes give \(\mu_{aj}\in[0,1]\), hence \(\tau_j\leq1\). % lean: CausalSmith.Stat.DiscreteBudgetvalueCurve.poRegression_unit_interval_budget

For an occupied cell, \cref{obj:ass:consistency,obj:ass:conditional-exchangeability,obj:ass:overlap} give \(s_a(q_j)=P(X=j,A=a)\geq\epsilon p_j>0\) and
\[
g_a(q_j)
=p_j\frac{P(X=j,A=a,Y=1)}{P(X=j,A=a)}
=p_j\mu_{aj}.
\]
For a null cell, \(q_j=0\) and the same identity has both sides zero. Substituting these identities into \cref{obj:def:dual-process}, and using \(p_j\geq0\), yields, for \(0\leq\lambda\leq1\),
\[
f_\lambda(q_j)=p_j\mu_{0j}+p_j[\tau_j-\lambda]_+,
\qquad
F_P(\lambda)=B(P)+\sum_{j=1}^d p_j[\tau_j-\lambda]_+.
\] % lean: CausalSmith.Stat.DiscreteBudgetvalueCurve.budget_dualProcess_identified

Fix \(b\in[0,1]\). If \(\pi\in\Pi_b(P)\), then \(0\leq\pi_j\leq1\) and \(\sum_jp_j\pi_j\leq b\) by \cref{obj:def:budget-policy-class}. For \(\lambda\in[0,1]\), the inequality
\(p_j\pi_j(\tau_j-\lambda)\leq p_j[\tau_j-\lambda]_+\)
holds whether \(\tau_j-\lambda\) is nonnegative or negative. Since \(\lambda\geq0\),
\[
\sum_{j=1}^d p_j\pi_j\tau_j
=\lambda\sum_{j=1}^d p_j\pi_j
 +\sum_{j=1}^d p_j\pi_j(\tau_j-\lambda)
\leq b\lambda+\sum_{j=1}^d p_j[\tau_j-\lambda]_+.
\] % lean: CausalSmith.Stat.DiscreteBudgetvalueCurve.weighted_budget_weak_duality
Taking the maximum over feasible policies and using \cref{obj:def:budget-value} gives
\[
V_b(P)\leq b\lambda+F_P(\lambda)
\quad(0\leq\lambda\leq1),
\qquad
V_b(P)\leq\inf_{0\leq\lambda\leq1}\{b\lambda+F_P(\lambda)\}.
\] % lean: CausalSmith.Stat.DiscreteBudgetvalueCurve.budget_dual_weak_bound

For the reverse inequality, the feasible policy set is a compact subset of \([0,1]^d\), contains the zero policy, and has a continuous gain. Thus a feasible maximizer \(\pi^\star\) exists. Applied to the affine budget and box constraints, the polyhedral normal-cone formula gives nonnegative multipliers \(\lambda^\star,\eta_j^-,\eta_j^+\) satisfying
\[
\begin{gathered}
\sum_jp_j\pi^\star_j\leq b,\qquad
\lambda^\star\Bigl(b-\sum_jp_j\pi^\star_j\Bigr)=0,\\
\eta_j^-\pi^\star_j=0,\qquad
\eta_j^+(1-\pi^\star_j)=0,\qquad
p_j(\lambda^\star-\tau_j)-\eta_j^-+\eta_j^+=0.
\end{gathered}
\] % lean: Causalean.Mathlib.Optimization.finite_knapsack_kkt
These conditions give the hinge identity cell by cell. It is immediate when \(p_j=0\). When \(p_j>0\), a coordinate with \(\pi^\star_j=0\) has \(\eta_j^+=0\) and \(\tau_j\leq\lambda^\star\); one with \(\pi^\star_j=1\) has \(\eta_j^-=0\) and \(\tau_j\geq\lambda^\star\); an interior coordinate has both box multipliers zero and \(\tau_j=\lambda^\star\). Consequently,
\[
p_j\pi^\star_j\tau_j
=\lambda^\star p_j\pi^\star_j
 +p_j[\tau_j-\lambda^\star]_+,
\qquad
\sum_jp_j\pi^\star_j\tau_j
=b\lambda^\star+\sum_jp_j[\tau_j-\lambda^\star]_+,
\]
where the second equality uses budget complementary slackness. % lean: Causalean.Mathlib.Optimization.finite_knapsack_dual_witness

If \(\lambda^\star\leq1\), this is a witness with price in \([0,1]\). If \(\lambda^\star>1\), then \(\tau_j\leq1\) makes every hinge at \(\lambda^\star\) zero. Also,
\[
b\lambda^\star
=\sum_jp_j\pi^\star_j\tau_j
\leq\sum_jp_j\pi^\star_j
\leq b.
\]
As \(b\geq0\), this forces \(b=0\). Every hinge at price \(1\) is then zero as well, and the displayed witness equality holds with price \(1\). Hence in every case there is \(\lambda_b\in[0,1]\) such that
\[
b\lambda_b+F_P(\lambda_b)
=B(P)+\sum_jp_j\pi^\star_j\tau_j
=V_b(P).
\] % lean: Causalean.Mathlib.Optimization.finite_knapsack_dual_witness
The process \(F_P\) is a finite sum of continuous hinge functions on \([0,1]\), so it is bounded there. Thus the dual infimum is well defined and is at most its value at \(\lambda_b\). Together with weak duality, this proves the representation. % lean: CausalSmith.Stat.DiscreteBudgetvalueCurve.dual_representation

Finally, define
\[
\omega_{\widehat F,F_P}=\sup_{\lambda\in[0,1]}|\widehat F(\lambda)-F_P(\lambda)|.
\]
Both processes are bounded on \([0,1]\), so \(\omega_{\widehat F,F_P}\) is finite and both dual objectives are bounded below for \(b\in[0,1]\). For every such \(b\) and every \(\lambda\in[0,1]\), the two inequalities
\[
b\lambda+\widehat F(\lambda)\leq b\lambda+F_P(\lambda)+\omega_{\widehat F,F_P},
\qquad
b\lambda+F_P(\lambda)\leq b\lambda+\widehat F(\lambda)+\omega_{\widehat F,F_P}
\]
give, on taking infima,
\[
\left|
\inf_{0\leq\lambda\leq1}\{b\lambda+\widehat F(\lambda)\}
-\inf_{0\leq\lambda\leq1}\{b\lambda+F_P(\lambda)\}
\right|\leq \omega_{\widehat F,F_P}.
\] % lean: CausalSmith.Stat.DiscreteBudgetvalueCurve.dualMinimum_uniform_contraction
Substitute the representation just proved and take the supremum over \(b\in[0,1]\). % lean: CausalSmith.Stat.DiscreteBudgetvalueCurve.dual_representation
\end{proof}
